\section{Conclusions}
\label{sec:discussion}

The competition posed a practical investment problem: pursue benchmark
outperformance through a systematic long-short portfolio while making its
risks and costs explicit. We addressed it with a framework connecting
historical data preparation, comparable strategy searches and shared portfolio
accounting. Trade audits and reconciled attribution connect the reported
returns to the underlying decisions and account movements.

Our return-focused comparison selected momentum \#1292 over more defensive
alternatives, available both within momentum and through sector momentum.
We evaluated the selected configuration on the competition window, where it
outperformed the benchmark after modeled costs. Interaction between sector
weights and relative stock returns was the largest attribution effect, led by
the concentrated long Information Technology sleeve. The short window and statistically insignificant
CAPM alpha in both the test and competition windows limit the evidence for
persistent risk-adjusted outperformance. The frontier identifies reallocations
with lower estimated volatility at the same estimated return, based on twelve
monthly observations. We have yet to evaluate a trading strategy using those
optimized weights.

Our next step would be to extend evaluation of the fixed strategy beyond
the competition window, comparing it with prespecified defensive alternatives
under wider spread and borrowing-cost assumptions. We would also test whether
rolling frontier allocations improve net performance. To reduce reliance on
the short sample, we would pull extreme estimated returns toward a common
average and investigate methods that retain older observations for stocks with
longer histories.
